% Fichier généré automatiquement par tools/generate_corrections.py.
% Source : SLCI/SLCI-03-SchemaBlocs/500_Divers/500_Divers.tex
% Statut : complet (4/4 question(s) renseignée(s), 0 reprise(s)).
\begin{corrigebox}[Corrigé extrait de la banque]
\CorrigeQuestion{1}
On a $\text{FTBO}(p)=\dfrac{K^2}{\left(R+Lp\right)\left(f+Jp\right)}$
$=\dfrac{K^2}{Rf+RJp+Lfp+LJp^2}$
$=\dfrac{K^2}{Rf\left(1+p\dfrac{RJ+Lf}{Rf}+\dfrac{LJ}{Rf}p^2\right)}$.

On a donc $K_{\text{BO}}=\dfrac{K^2}{Rf}$, 
$\omega_{\text{BO}} = \sqrt{\dfrac{Rf}{LJ}}$,
$\dfrac{2\xi_{\text{BO}} }{\omega_{\text{BO}}}=\dfrac{RJ+Lf}{Rf} \Leftrightarrow
\xi_{\text{BO}} =\omega_{\text{BO}}\dfrac{RJ+Lf}{2Rf}
=\sqrt{\dfrac{Rf}{LJ}}\dfrac{RJ+Lf}{2Rf}
=\dfrac{RJ+Lf}{2\sqrt{LJRf}}$.
\CorrigeQuestion{2}
On a $\text{FTBF}(p)=\dfrac{\dfrac{K}{\left(R+Lp\right)\left(f+Jp\right)}}{1+\dfrac{K^2}{\left(R+Lp\right)\left(f+Jp\right)}}
=\dfrac{K}{\left(R+Lp\right)\left(f+Jp\right)+K^2}
=\dfrac{\dfrac{K}{K^2+Rf}}{\dfrac{RJ+Lf}{Rf+K^2}p+\dfrac{LJ}{Rf+K^2}p^2+1}$.

On a donc $K_{\text{BF}}=\dfrac{K}{K^2+Rf}$, 
$\omega_{\text{BF}} = \sqrt{\dfrac{Rf+K^2}{LJ}}$,
$\dfrac{2\xi_{\text{BF}} }{\omega_{\text{BF}}}=\dfrac{RJ+Lf}{Rf+K^2} \Leftrightarrow
\xi_{\text{BO}} =\omega_{\text{BF}}\dfrac{RJ+Lf}{2\left(Rf+K^2\right)}
=\sqrt{\dfrac{Rf+K^2}{LJ}}\dfrac{RJ+Lf}{2\left(Rf+K^2\right)}$

$\xi_{\text{BF}}=\dfrac{RJ+Lf}{2\sqrt{LJ}\sqrt{Rf+K^2}}$.
\CorrigeQuestion{3}
Si on note $R(p)$ la seconde entrée du \textbf{premier comparateur} et $\varepsilon(p)$ la sortie du premier comparateur,  

$\text{FTBO(p)}=\dfrac{\varepsilon(p)}{R(p)} = A \times \dfrac{\dfrac{1}{p}}{1+\dfrac{B}{p}}\times C = \dfrac{AC}{B+p} = \dfrac{\dfrac{AC}{B}}{1+\dfrac{p}{B}}$.
On a donc $K_{\text{BO}}=\dfrac{AC}{B}$ et $\tau_{\text{BO}}=\dfrac{1}{B}$.
\CorrigeQuestion{4}
On a
$\text{FTBF(p)} = \dfrac{\dfrac{A}{B+p}}{1+\dfrac{AC}{B+p}}=\dfrac{A}{B+p+AC}=\dfrac{\dfrac{A}{B+AC}}{1+\dfrac{p}{B+AC}}$.

On a donc $K_{\text{BF}}=\dfrac{A}{B+AC}$ et $\tau_{\text{BF}}=\dfrac{1}{B+AC}$.

\begin{enumerate}
\item $K_{\text{BO}}=\dfrac{K^2}{Rf}$, 
$\omega_{\text{BO}} = \sqrt{\dfrac{Rf}{LJ}}$,
$\xi_{\text{BO}} =\dfrac{RJ+Lf}{2\sqrt{LJRf}}$.
\item $K_{\text{BF}}=\dfrac{K}{K^2+Rf}$, 
$\xi_{\text{BF}}=\dfrac{RJ+Lf}{2\sqrt{LJ}\sqrt{Rf+K^2}}$.
\item $K_{\text{BO}}=\dfrac{AC}{B}$ et $\tau_{\text{BO}}=\dfrac{1}{B}$.
\item $K_{\text{BF}}=\dfrac{A}{B+AC}$ et $\tau_{\text{BF}}=\dfrac{1}{B+AC}$.
\end{enumerate}
\end{corrigebox}
